% Volume V: Dynamics, Signals, Graphs, and Control
% Continuity note for Chapter 27.
% Previous dependencies: Chapters 13--15 provide probability spaces, random variables, conditional distributions, expectation, and Bayes updates; Chapters 25--26 provide deterministic dynamics and differential-equation notation.
% New durable concepts: stochastic processes as indexed random variables; sample paths; finite-dimensional distributions; discrete-time transition kernels; Markov chains; transition matrices; stationary distributions; absorbing and transient states; Brownian motion; Ito SDE notation; diffusion scaling; Ornstein--Uhlenbeck processes; Poisson counting processes; intensities and hazards; sequential Bayes filters; Kalman-filter and particle-filter mechanisms.
% Forward hooks: Chapter 28 uses autocovariance and spectral summaries of time series; Chapter 30 uses Markov state dynamics in dynamic programming; Chapter 36 turns Markov chains into Markov decision processes; Chapter 41 uses SDE intuition for diffusion models; Chapter 46 specializes point processes to spike trains and neural GLMs.
% Notation commitments: Stochastic processes are X_t or X_n on a probability space (\Omega,\mathcal F,P); sample paths are t\mapsto X_t(\omega); finite-dimensional distributions are laws of (X_{t_1},\ldots,X_{t_k}); row distributions for finite Markov chains use \pi_{n+1}=\pi_n P; Brownian motion is W_t; Ito SDEs use dX_t=a(X_t,t)dt+b(X_t,t)dW_t.
\section{Chapter 27: Stochastic Processes, Markov Chains, and SDEs}
\label{ch:stochastic-processes-markov-chains-sdes}

Deterministic dynamics ask for a trajectory once the initial condition and the vector field are known. Neural and learning systems often require a different object. Spikes fluctuate from trial to trial, synaptic input contains random arrivals, behavioral choices vary, minibatches perturb gradients, and latent states must be inferred from noisy observations. The central mathematical object is no longer one curve but a probability law over curves.

\paragraph{Chapter problem.}
How do we represent random time-evolving quantities, compute with discrete and continuous stochastic dynamics, and update beliefs about hidden states as observations arrive?

\paragraph{Section map.}
Section 27.1 defines stochastic processes, sample paths, finite-dimensional distributions, and discrete-time transition kernels. Section 27.2 develops Markov chains through transition matrices, stationary distributions, and absorbing/transient behavior. Section 27.3 introduces continuous-time processes, Brownian motion, SDEs, Ornstein--Uhlenbeck noise, and Poisson counting processes. Section 27.4 turns stochastic dynamics into filtering and state estimation through sequential Bayes, Kalman filtering, and particle filtering.

\subsection{27.1 Discrete-time stochastic processes}

\paragraph{Guiding question.}
What is the mathematical object behind a random time series such as a sequence of spike counts, rewards, or noisy gradients?

\paragraph{Beginner-facing intuition.}
A deterministic sequence is a list of numbers \(x_0,x_1,x_2,\ldots\). A stochastic process is a list of random variables
\[
X_0,X_1,X_2,\ldots.
\]
Before the experiment, each \(X_n\) is uncertain. After one outcome \(\omega\) occurs, the realized values
\[
X_0(\omega),X_1(\omega),X_2(\omega),\ldots
\]
form one sample path. A neural experiment may produce many trial-by-trial sample paths of a spike-count process. A training run may produce one sample path of a noisy loss process.

The point of the notation is to separate three layers: the probability model, the random variables indexed by time, and one realized path. Confusing those layers is the most common beginner error in stochastic modeling.

\paragraph{Formal definitions and notation.}
Let \((\Omega,\mathcal F,P)\) be a probability space and let \(T\) be an index set. A \emph{stochastic process} with state space \(S\) is a family of random variables
\[
\{X_t:t\in T\},\qquad X_t:\Omega\to S.
\]
If \(T=\{0,1,2,\ldots\}\), the process is discrete-time. If \(T=[0,\infty)\), it is continuous-time. For a fixed outcome \(\omega\), the map
\[
t\mapsto X_t(\omega)
\]
is a \emph{sample path}. For times \(t_1,\ldots,t_k\), the joint law of
\[
(X_{t_1},\ldots,X_{t_k})
\]
is a \emph{finite-dimensional distribution}. A process is specified, in distribution, by all of these joint laws when the usual consistency conditions hold.

A discrete-time process with state space \(S\) has a \emph{transition kernel} when, for each current state \(x\in S\), there is a probability distribution
\[
K(x,B)=P(X_{n+1}\in B\mid X_n=x)
\]
over measurable sets \(B\subseteq S\). If \(S\) is finite, the kernel is a matrix.

\paragraph{Core mechanism: random walk moments.}
Let \(\xi_1,\xi_2,\ldots\) be independent random variables with
\[
P(\xi_i=1)=p,\qquad P(\xi_i=-1)=1-p.
\]
Define the random walk
\[
X_0=0,\qquad X_n=\sum_{i=1}^n \xi_i.
\]
Then
\[
\mathbb E[X_n]=n(2p-1),
\qquad
\operatorname{Var}(X_n)=4np(1-p).
\]
\emph{Derivation.} Each increment has
\[
\mathbb E[\xi_i]=p-(1-p)=2p-1
\]
and, since \(\xi_i^2=1\),
\[
\operatorname{Var}(\xi_i)=1-(2p-1)^2=4p(1-p).
\]
Linearity of expectation gives the mean. Independence makes variances add, giving the variance. For \(p=1/2\), the expected position is \(0\), but the standard deviation is \(\sqrt n\), so typical paths wander farther from the origin as time grows.

\paragraph{Concrete example.}
For \(p=1/2\) and \(n=100\),
\[
\mathbb E[X_{100}]=0,\qquad \operatorname{Var}(X_{100})=100,\qquad \operatorname{SD}(X_{100})=10.
\]
One realized path may end at \(14\), another at \(-6\), and another at \(0\). None of those individual endings is the expectation. The expectation summarizes the ensemble.

\paragraph{Computational neuroscience example.}
Let \(X_n\) be the binned spike count of one neuron in the \(n\)-th \(20\) ms bin of a trial. Across repeated trials, \(X_n\) is a random variable. A single plotted raster or binned trace is a sample path. The joint distribution of \((X_{10},X_{11},X_{12})\) describes temporal dependence across nearby bins. If refractoriness or slow gain changes are present, this joint law is not determined by separate marginal histograms at each bin.

\paragraph{AI or machine-learning example.}
In stochastic gradient descent, let \(\Theta_n\) be the parameter vector after \(n\) minibatch updates. Then \(\{\Theta_n\}\) is a stochastic process because the minibatch sequence is random. A loss curve from one training run is one sample path of \(L(\Theta_n)\). Averaging over many random seeds estimates properties of the process, not a deterministic training trajectory.

\paragraph{Common misconceptions and failure modes.}
A random variable \(X_n\) is not the same thing as an observed value \(x_n\). A sample path is not the whole process; it is one realization from the process. Matching all one-time marginals \(P(X_n\in B)\) is not enough to match the process, because temporal dependence lives in joint distributions. Finally, a transition kernel describes one-step conditional behavior; it is not automatically valid if the next state depends on deeper history.

\paragraph{Exercises.}
\begin{enumerate}[label=\textbf{27.1.\arabic*.}, leftmargin=*]
\item \textbf{Mechanical.} Let \(X_n\) be a spike count in bin \(n\). Identify the probability space, the index set, the state space, and one possible sample path in words.
\item \textbf{Proof or derivation.} For the random walk above, derive \(\mathbb E[X_n]\) and \(\operatorname{Var}(X_n)\) using linearity and independence.
\item \textbf{Numerical/computational.} For \(p=0.6\), compute \(\mathbb E[X_{50}]\) and \(\operatorname{SD}(X_{50})\).
\item \textbf{Interpretation.} Two processes have the same distribution for every single \(X_n\), but one alternates signs while the other has independent signs. Why do their finite-dimensional distributions differ?
\item \textbf{Misconception/debugging.} A student simulates one path and says, ``The process has mean \(7\) at time \(100\).'' What additional information is needed before that statement is justified?
\end{enumerate}

\subsection{27.2 Markov chains}

\paragraph{Guiding question.}
When is the current state enough information to predict the distribution of the next state?

\paragraph{Beginner-facing intuition.}
A Markov chain is a stochastic process with a short memory. Once the present state is known, the past does not add information about the next step. This does not mean the process is independent across time. It means the dependence is carried through the current state.

This is a good first model for neural or behavioral state switching: sleep stages, latent decision states, active/inactive circuit states, or token-by-token generation with a finite context state.

\paragraph{Formal definitions and notation.}
Let \(S=\{1,\ldots,m\}\) be a finite state space. A discrete-time process \(\{X_n\}_{n\geq0}\) is a \emph{time-homogeneous Markov chain} if
\[
P(X_{n+1}=j\mid X_n=i,X_{n-1}=i_{n-1},\ldots,X_0=i_0)
=P(X_{n+1}=j\mid X_n=i)
\]
whenever the conditioning event has positive probability, and the right side does not depend on \(n\). The transition matrix \(P\in\mathbb R^{m\times m}\) has entries
\[
P_{ij}=P(X_{n+1}=j\mid X_n=i),
\qquad
P_{ij}\geq0,\qquad \sum_{j=1}^m P_{ij}=1.
\]
If \(\pi_n\) is the row-vector distribution of \(X_n\), then
\[
\boxed{\pi_{n+1}=\pi_n P,\qquad \pi_n=\pi_0P^n.}
\]
A distribution \(\pi_\star\) is \emph{stationary} if
\[
\pi_\star=\pi_\star P,\qquad \sum_i \pi_{\star,i}=1,\qquad \pi_{\star,i}\geq0.
\]
A state \(i\) is \emph{absorbing} if \(P_{ii}=1\). A state is \emph{transient} if, starting from that state, there is positive probability of leaving and never returning.

\paragraph{Core mechanism: stationary distribution from flow balance.}
Consider the two-state chain
\[
P=
\begin{pmatrix}
1-a & a\\
b & 1-b
\end{pmatrix},
\qquad 0<a,b<1.
\]
Let \(\pi_\star=(\pi_1,\pi_2)\). The equation \(\pi_\star=\pi_\star P\) gives
\[
\pi_1 a=\pi_2 b,\qquad \pi_1+\pi_2=1.
\]
Thus
\[
\boxed{\pi_\star=\left(\frac{b}{a+b},\frac{a}{a+b}\right).}
\]
The first equation says that, at stationarity, probability flow from state \(1\) to state \(2\) equals probability flow from state \(2\) to state \(1\). This flow-balance intuition generalizes, though larger chains require solving a linear system.

\paragraph{Concrete example.}
Let
\[
P=
\begin{pmatrix}
0.8 & 0.2\\
0.1 & 0.9
\end{pmatrix}.
\]
If \(\pi_0=(1,0)\), then
\[
\pi_1=(0.8,0.2)
\]
and
\[
\pi_2=\pi_1P=(0.8,0.2)
\begin{pmatrix}
0.8 & 0.2\\
0.1 & 0.9
\end{pmatrix}
=(0.66,0.34).
\]
The stationary distribution is
\[
\pi_\star=\left(\frac{0.1}{0.3},\frac{0.2}{0.3}\right)=\left(\frac13,\frac23\right).
\]
The second state is more common at stationarity because the chain leaves state \(2\) slowly.

\paragraph{Computational neuroscience example.}
Suppose a population switches between a low-activity state \(L\) and a high-activity state \(H\). If
\[
P(L\to H)=0.02,\qquad P(H\to L)=0.05
\]
per time bin, then the stationary high-state probability is
\[
\frac{0.02}{0.02+0.05}\approx0.286.
\]
This is not a claim that cortical dynamics are exactly two-state Markov. It is a controlled summary: transition probabilities determine occupancy and dwell-time statistics under the Markov approximation.

\paragraph{AI or machine-learning example.}
A language model with a finite hidden state can be idealized as a Markov chain over hidden states or tokens. In reinforcement learning, a Markov decision process adds actions to this idea: the next-state distribution depends on the current state and chosen action. Chapter 30 uses this extension for Bellman equations, and Chapter 36 uses it for value and policy learning.

\paragraph{Common misconceptions and failure modes.}
The Markov property does not mean the states are independent. Adjacent states can be strongly dependent. A stationary distribution is not necessarily the initial distribution; it is a distribution unchanged by one transition. An absorbing state can dominate long-run behavior even if it is rarely entered. A finite transition matrix must have rows summing to one under the row-vector convention used here; using a column convention without saying so reverses update formulas.

\paragraph{Exercises.}
\begin{enumerate}[label=\textbf{27.2.\arabic*.}, leftmargin=*]
\item \textbf{Mechanical.} Check whether
\[
P=\begin{pmatrix}0.7&0.3\\0.4&0.6\end{pmatrix}
\]
is a valid row-stochastic transition matrix.
\item \textbf{Proof or derivation.} Derive the stationary distribution of the two-state chain with transition probabilities \(a\) and \(b\) as in the section.
\item \textbf{Numerical/computational.} For \(\pi_0=(0,1)\) and \(P=\begin{pmatrix}0.8&0.2\\0.1&0.9\end{pmatrix}\), compute \(\pi_1\) and \(\pi_2\).
\item \textbf{Interpretation.} A state has \(P_{ii}=1\). What does that imply for future trajectories once the chain enters that state?
\item \textbf{Misconception/debugging.} A student says, ``Because the process is Markov, \(X_{n+1}\) is independent of \(X_n\).'' Correct the statement.
\end{enumerate}

\subsection{27.3 Continuous-time processes and SDEs}

\paragraph{Guiding question.}
How can a trajectory be continuous in time while still being driven by random fluctuations or random events?

\paragraph{Beginner-facing intuition.}
Continuous-time stochastic models come in two common forms. Some processes move continuously but irregularly, like Brownian motion or noisy membrane voltage. Others jump by discrete events in continuous time, like spike counts or synaptic arrivals. SDEs model the first kind by adding noise to differential equations. Point processes model the second kind by random event times.

The notation \(dW_t\) in an SDE is not an ordinary tiny deterministic number. It represents Brownian fluctuation whose size is on the order of \(\sqrt{dt}\), not \(dt\). This scaling is why stochastic calculus is not just ordinary calculus with a random term added.

\paragraph{Formal definitions and notation.}
A standard \emph{Brownian motion}, or Wiener process, is a continuous-time process \(\{W_t:t\geq0\}\) such that
\[
W_0=0,\qquad W_t-W_s\sim\mathcal N(0,t-s)\quad(0\leq s<t),
\]
increments over disjoint intervals are independent, and sample paths are continuous with probability one. Brownian paths are almost surely nowhere differentiable, even though they are continuous.

An Ito SDE in one dimension has the form
\[
\boxed{dX_t=a(X_t,t)\,dt+b(X_t,t)\,dW_t.}
\]
The function \(a\) is the drift and \(b\) is the diffusion amplitude. Over a small step \(\Delta t\), Euler--Maruyama simulation uses
\[
X_{n+1}=X_n+a(X_n,t_n)\Delta t+b(X_n,t_n)\sqrt{\Delta t}\,Z_n,
\qquad Z_n\sim\mathcal N(0,1).
\]

A \emph{counting process} \(N(t)\) records the number of events up to time \(t\). A homogeneous Poisson process with rate \(\lambda>0\) satisfies
\[
N(t)-N(s)\sim\mathrm{Pois}(\lambda(t-s))
\]
for \(0\leq s<t\), and increments over disjoint intervals are independent. An inhomogeneous Poisson process has deterministic intensity \(\lambda(t)\), with
\[
N(a,b)\sim \mathrm{Pois}\!\left(\int_a^b \lambda(u)\,du\right).
\]
For a general point process, the conditional intensity or hazard \(\lambda(t\mid\mathcal H_t)\) satisfies
\[
P\{N(t,t+dt)=1\mid\mathcal H_t\}
=\lambda(t\mid\mathcal H_t)\,dt+o(dt),
\]
where \(\mathcal H_t\) is the history before \(t\).

\paragraph{Core mechanisms: diffusion scaling, OU noise, and Poisson waiting times.}
\emph{Brownian scaling from random walks.} Let
\[
B^{(\Delta t)}(t)=\sqrt{\Delta t}\sum_{i=1}^{\lfloor t/\Delta t\rfloor}\xi_i,
\]
where \(\xi_i=\pm1\) with equal probability. Then
\[
\mathbb E[B^{(\Delta t)}(t)]=0,\qquad
\operatorname{Var}(B^{(\Delta t)}(t))\approx t.
\]
The \(\sqrt{\Delta t}\) step size is the only scaling that keeps variance finite and nonzero as \(\Delta t\to0\). Brownian motion is the continuous-time diffusion limit of such rescaled random walks.

\emph{Ornstein--Uhlenbeck process.} The OU process is
\[
dX_t=\theta(\mu-X_t)\,dt+\sigma\,dW_t,\qquad \theta>0.
\]
The drift pulls \(X_t\) back toward \(\mu\). Its conditional mean is
\[
\mathbb E[X_t\mid X_0=x_0]=\mu+(x_0-\mu)e^{-\theta t}.
\]
At stationarity,
\[
X_t\sim\mathcal N\!\left(\mu,\frac{\sigma^2}{2\theta}\right),
\qquad
\operatorname{Cov}(X_t,X_{t+\tau})=\frac{\sigma^2}{2\theta}e^{-\theta|\tau|}.
\]
Thus \(\theta^{-1}\) is the correlation time scale.

\emph{Poisson interarrival times.} For a homogeneous Poisson process,
\[
P(T_1>t)=P(N(t)=0)=e^{-\lambda t}.
\]
Therefore the first arrival time has density
\[
f_{T_1}(t)=\lambda e^{-\lambda t},\qquad t\geq0.
\]
The constant hazard \(\lambda\) is the memoryless event-rate property.

\paragraph{Concrete examples.}
For an OU process with \(\theta=2\), \(\mu=0\), and \(\sigma=1\), the stationary variance is \(1/4\), and the autocovariance is
\[
C(\tau)=\frac14 e^{-2|\tau|}.
\]
For a homogeneous Poisson process with \(\lambda=20\) spikes/s, the count in a \(0.1\) s window has mean and variance
\[
\lambda T=20(0.1)=2.
\]
The probability of no spikes in that window is \(e^{-2}\).

\paragraph{Computational neuroscience example.}
OU noise is a common model for fluctuating synaptic drive or latent gain because it is continuous, mean-reverting, Gaussian, and temporally correlated. A membrane-potential approximation might write
\[
dV_t=-\frac{1}{\tau}(V_t-V_{\mathrm{rest}})\,dt+\sigma\,dW_t,
\]
before adding threshold and reset. Poisson processes model idealized spike arrivals or presynaptic event counts. The Poisson assumption is a baseline, not a claim that real neurons lack refractoriness, bursting, adaptation, or shared drive.

\paragraph{AI or machine-learning example.}
Diffusion and score-based generative models use SDE-like ideas: a forward process gradually adds noise, and a learned reverse process tries to remove it. Stochastic optimization can also be approximated by an SDE when small learning rates and minibatch noise produce many small random parameter increments. In both cases, the Brownian scaling says why accumulated small noise has variance proportional to time.

\paragraph{Common misconceptions and failure modes.}
The expression \(dW_t/dt\) is not an ordinary derivative; Brownian paths are too rough. The noise increment \(dW_t\) has size \(\sqrt{dt}\), so it cannot be treated as the same order as \(dt\). A Poisson process is not any count process; it requires independent increments and a specified intensity. A time-varying rate in an inhomogeneous Poisson process creates nonstationary counts but still leaves independent increments conditional on the rate function.

\paragraph{Exercises.}
\begin{enumerate}[label=\textbf{27.3.\arabic*.}, leftmargin=*]
\item \textbf{Mechanical.} State the mean and variance of \(W_t-W_s\) for Brownian motion.
\item \textbf{Proof or derivation.} Derive the stationary variance \(\sigma^2/(2\theta)\) of the OU process by using the known Gaussian stationary law or by balancing mean reversion and diffusion at the covariance level.
\item \textbf{Numerical/computational.} Write the Euler--Maruyama update for \(dX_t=-3X_t\,dt+0.5\,dW_t\) with step size \(\Delta t=0.01\).
\item \textbf{Interpretation.} A spike-count process has \(P(N(0.1)=0)=e^{-5}\) for a \(0.1\) s window under a homogeneous Poisson model. What is the firing rate \(\lambda\)?
\item \textbf{Misconception/debugging.} A student says, ``An SDE is just an ODE where \(dW_t\) is a small deterministic input.'' Explain why the scaling and path regularity make this wrong.
\end{enumerate}

\subsection{27.4 Filtering and state estimation}

\paragraph{Guiding question.}
How can a model combine noisy dynamics with noisy observations to estimate a hidden state over time?

\paragraph{Beginner-facing intuition.}
Filtering is sequential Bayesian inference. A hidden state evolves through time, observations arrive one by one, and the posterior after one observation becomes the prior information for the next step. The loop is always the same:
\[
\text{predict using dynamics, then update using the new observation.}
\]
This is the mathematical skeleton behind neural decoding, tracking, robotics, state-space models, and many recurrent inference algorithms.

\paragraph{Formal definitions and notation.}
A discrete-time state-space model has hidden states \(X_t\in S\) and observations \(Y_t\in O\). A common model specifies
\[
p(x_t\mid x_{t-1})\qquad\text{and}\qquad p(y_t\mid x_t).
\]
The filtering distribution is
\[
p(x_t\mid y_{1:t}),
\qquad y_{1:t}=(y_1,\ldots,y_t).
\]
The Bayes filter recursion is
\[
\boxed{p(x_t\mid y_{1:t-1})
=\int p(x_t\mid x_{t-1})p(x_{t-1}\mid y_{1:t-1})\,dx_{t-1}}
\]
followed by
\[
\boxed{p(x_t\mid y_{1:t})
=\frac{p(y_t\mid x_t)p(x_t\mid y_{1:t-1})}
{\int p(y_t\mid x_t)p(x_t\mid y_{1:t-1})\,dx_t}.}
\]
For finite states, the integrals become sums.

\paragraph{Core mechanism: scalar Kalman update.}
In the linear-Gaussian model
\[
X_t=aX_{t-1}+\epsilon_t,\qquad \epsilon_t\sim\mathcal N(0,q),
\]
\[
Y_t=cX_t+\eta_t,\qquad \eta_t\sim\mathcal N(0,r),
\]
Gaussian beliefs remain Gaussian. If
\[
X_{t-1}\mid y_{1:t-1}\sim\mathcal N(m_{t-1},P_{t-1}),
\]
then the prediction step is
\[
m_t^- = a m_{t-1},\qquad P_t^- = a^2P_{t-1}+q.
\]
After observing \(y_t\), define the Kalman gain
\[
K_t=\frac{P_t^- c}{c^2P_t^-+r}.
\]
The update is
\[
m_t=m_t^-+K_t(y_t-cm_t^-),
\qquad
P_t=(1-K_tc)P_t^-.
\]
The term \(y_t-cm_t^-\) is the innovation: the part of the observation not predicted by the current belief.

\paragraph{Concrete example.}
Let \(a=c=1\), \(q=1\), \(r=4\), \(m_{t-1}=0\), and \(P_{t-1}=1\). The prediction gives
\[
m_t^-=0,\qquad P_t^-=2.
\]
The Kalman gain is
\[
K_t=\frac{2}{2+4}=\frac13.
\]
If \(y_t=6\), then
\[
m_t=0+\frac13(6)=2,
\qquad
P_t=(1-\frac13)2=\frac43.
\]
The posterior mean moves toward the observation but does not equal it, because measurement noise is substantial.

\paragraph{Computational neuroscience example.}
In neural decoding, \(X_t\) may be a hand position, head direction, latent decision variable, or arousal state, while \(Y_t\) is a population spike-count vector. The dynamics model predicts how the latent state moves between time bins. The observation model says how neural activity depends on the state. Filtering combines both: smooth latent trajectories are not inferred from spikes alone, and neural observations are not interpreted without a dynamical prior.

\paragraph{AI or machine-learning example.}
Particle filters approximate the filtering distribution by weighted samples
\[
\{x_t^{(i)},w_t^{(i)}\}_{i=1}^M.
\]
Each particle is propagated through the dynamics, reweighted by the likelihood of the new observation, and sometimes resampled. This is a Monte Carlo version of the Bayes filter. It is useful when the state-space model is nonlinear or non-Gaussian, where exact Kalman formulas do not apply.

\paragraph{Common misconceptions and failure modes.}
Filtering is not smoothing: filtering uses observations up to time \(t\), while smoothing may use future observations. A Kalman filter is not a generic synonym for any filter; it relies on linear-Gaussian assumptions or linearized approximations. A particle filter does not remove approximation error automatically; too few particles can collapse onto a small number of hypotheses. Finally, the update step is not just averaging the prior and observation. The weights are set by uncertainty and likelihood.

\paragraph{Exercises.}
\begin{enumerate}[label=\textbf{27.4.\arabic*.}, leftmargin=*]
\item \textbf{Mechanical.} Name the two distributions needed to define a state-space model: the state transition model and the observation model.
\item \textbf{Proof or derivation.} Derive the finite-state Bayes filter update by applying Bayes' theorem to \(p(x_t\mid y_{1:t})\) after the prediction step.
\item \textbf{Numerical/computational.} In the scalar Kalman example, recompute the gain and posterior mean when \(r=1\) instead of \(r=4\). Interpret the change.
\item \textbf{Interpretation.} Why can filtering improve neural decoding when spike counts are noisy but the latent movement state evolves smoothly?
\item \textbf{Misconception/debugging.} A student says, ``After seeing \(y_t\), the best estimate of \(x_t\) is simply \(y_t\).'' Explain why this ignores the observation model, units, and measurement noise.
\end{enumerate}

\paragraph{Closing bridge.}
Stochastic processes give time-dependent uncertainty a precise shape. Markov chains make discrete state switching computable, SDEs and point processes model continuous-time noise and events, and filters turn noisy observations into evolving beliefs. The next chapter studies how time series are represented, decomposed, filtered, and summarized in frequency space; those tools are often applied to the very stochastic processes introduced here.

\clearpage
